\section{Integral cohomology and degree-three group classes}
\label{sec:cohomology}

Our goal is to construct the degree-three classes that will be transported
through the two-height tower, together with the exact restriction maps
between their linear realizations. We first fix the integral coefficient
convention and prove the finite-resolution statement that will also
control the later Galois-weight and ordinary-descent arguments.

\subsection{Coefficients and continuous descent}

Put $\Lambda_m=\Z/p^m$.  For a small v-stack $X$ occurring below, put
\[
 R\Gamma_{\mathrm{int}}(X)
   =\Rlim_m R\Gamma_v(X,\Lambda_m),\qquad
 H^i_{\mathrm{int}}(X)=H^iR\Gamma_{\mathrm{int}}(X),\qquad
 H^i_{\mathrm b}(X)=H^i_{\mathrm{int}}(X)[1/p].
\]
For a locally profinite group $G$, define separately in $D(\mathrm{Ab})$
\begin{equation}\label{eq:integral-group-cochains}
\begin{gathered}
 R\Gamma_{\mathrm{int}}(G)=\Rlim_m C_{\cts}^{\bullet}(G,\Lambda_m),
 \\
 H^i_{\mathrm{int}}(G)=H^iR\Gamma_{\mathrm{int}}(G),
 \qquad H^i_{\mathrm b}(G)=H^i_{\mathrm{int}}(G)[1/p].
\end{gathered}
\end{equation}
Here $C_{\cts}^{\bullet}$ is the inhomogeneous continuous cochain complex,
with trivial discrete coefficients.  The standard resolution in the
condensed classifying topos $BG_{\mathrm{pro\acute et}}$, the topos of
condensed sets with $G$-action, gives the natural comparison
\[
 C_{\cts}^{\bullet}(G,\Lambda_m)\simeq
 \RHom_{\Lambda_m[\underline G]\text{-}\mathrm{Cond}}
       (\underline{\Lambda_m},\underline{\Lambda_m}).
\]
This is the external derived Hom; its continuous-cochain calculation for
these solid topological coefficients is the standard-resolution proof in
\cite[\S2, pp.~5--7, Lemmas~2.1--2.2]{Anschutz2020}.
We next construct the maps from this group complex to geometric quotients.

For a condensed $G$-complex $K$, write $R\Gamma(G,K)$ for its
\emph{internal} derived condensed invariants, and $R\Gamma(G,K)(*)$
for their point value, which is the external derived Hom
\cite[\S2, p.~8]{Anschutz2020}.  Topological coefficients in this
notation are implicitly condensed.  Reduction on finite cochains is
degreewise surjective: a function lifts by composing with a section of
the finite coefficient reduction.  Compatible finite continuous
functions are exactly continuous $\Z_p$-valued functions.  Hence
\begin{equation}\label{eq:group-cochain-bridge}
 R\Gamma_{\mathrm{int}}(G)\simeq C_{\cts}^{\bullet}(G,\Z_p)
       \simeq R\Gamma(G,\underline{\Z_p})(*).
\end{equation}
For the last comparison, $\Z_p$ with its $p$-adic topology is solid by
\cite[\S1, pp.~4--5]{Anschutz2020}, so the same standard-resolution
proof applies.  We do not identify the internal invariant object with
the discrete condensation of the external cochain complex for a
noncompact group; see \cite[Lemma~2.5 and Remark~2.6]{Anschutz2020}.
The comparisons in \eqref{eq:group-cochain-bridge} are natural for
continuous group homomorphisms.
For compact $G$, continuous functions $G^r\to\Q_p$ are bounded, so
\[
 C_{\cts}(G^r,\Z_p)[1/p]=C_{\cts}(G^r,\Q_p).
\]
Consequently $H^*_{\mathrm b}(G)$ agrees with ordinary continuous
$\Q_p$-cohomology for compact $G$.  We use this last identification
only in that case.

Tate twists are taken at each finite level before taking the derived
limit.  All maps of coefficients and all descent maps are formed
integrally.  In particular, $H^i_{\mathrm b}(X)$ is not defined as
$H^i_v(X,\widehat\Z_p[1/p])$.  Moving inversion of $p$ before
cohomology can change the result for a noncompact space or group.

For an analytic adic space in this paper, its \'etale site agrees with
the \'etale site of its diamond
\cite[Lemma~15.6]{ScholzeDiamonds}. On a locally spatial diamond,
bounded-below \'etale complexes have the same derived cohomology on
the diamond \'etale, quasi-pro-\'etale, and v-sites
\cite[Proposition~14.10]{ScholzeDiamonds}. We apply these statements
to finite constant coefficients, one level at a time. The integral
analytic pro-\'etale comparison used for Drinfeld space is the
separate comparison in \cite[\S5]{CDN2021}. Quotient stacks are
handled by descent for their covers. Only after the finite-level
comparisons do we take $\Rlim_m$.

Here is the map from group classes to geometric quotients.  For a small
v-stack base $B$, put $B_BG=[B/\underline G]$, where $G$ acts trivially
on $B$.  A $G$-space $X$ over $B$ gives a map $[X/G]\to B_BG$; a
$G$-torsor $T\to Z$ over $B$ gives its classifying map $Z\to B_BG$.
At each finite level define
\begin{equation}\label{eq:finite-constant-cochains}
\begin{split}
 u_{X,G,m}:C_{\cts}^{\bullet}(G,\Lambda_m)
 &\longrightarrow
   \operatorname{Tot}_{a}R\Gamma_v
       (X\times\underline{G^a},\Lambda_m)\\
 &\simeq R\Gamma_v([X/G],\Lambda_m).
\end{split}
\end{equation}
In degree $a$, pull a locally constant function on $G^a$ back along
the projection from $X\times\underline{G^a}$.  Its finitely many
clopen fibers define a section of the finite constant sheaf, followed
by the canonical degree-zero truncation map to derived sections.
These maps commute with
the faces and degeneracies of the action groupoid.  The last
equivalence is its \v{C}ech descent.  The case $X=B$ defines the
universal map into $B_BG$, and the general map is its pullback along
$[X/G]\to B_BG$.  Likewise the map obtained from the \v{C}ech nerve
of $T\to Z$ is the pullback along its classifying map.

If $\phi:G'\to G$ is a continuous homomorphism and $f:X'\to X$ is
$\phi$-equivariant over a base map, this construction gives a
commutative square
\begin{equation}\label{eq:constant-cochain-naturality}
\begin{tikzcd}[column sep=large]
 C_{\cts}^{\bullet}(G,\Lambda_m)
     \arrow[r,"u_{X,G,m}"] \arrow[d,"\phi^*"']
   &R\Gamma_v([X/G],\Lambda_m)\arrow[d,"{[f/\phi]^*}"]\\
 C_{\cts}^{\bullet}(G',\Lambda_m)
     \arrow[r,"u_{X',G',m}"']
   &R\Gamma_v([X'/G'],\Lambda_m).
\end{tikzcd}
\end{equation}
It commutes already in each finite cosimplicial degree.  The same
naturality holds for coefficient reduction, maps of finite coefficient
sheaves, and automorphisms of the base.  Taking $\Rlim_m$ defines
$u_{X,G,\mathrm{int}}:R\Gamma_{\mathrm{int}}(G)\to
R\Gamma_{\mathrm{int}}([X/G])$; denote the induced maps on cohomology
after $[1/p]$ by $u_{X,G,\mathrm b}$.  No acyclicity of $B$ or $X$
is used in this construction.

There is one restricted case in which the universal map is an
equivalence.  Let $B_C=\Spd C$ for the algebraically closed perfectoid
field $C=\C_p$, and abbreviate $B_CG=[B_C/\underline G]$.
For every profinite set $S$, the actual constants map
is an equivalence
\begin{equation}\label{eq:geometric-point-constants}
 C_{\cts}(S,\Lambda_m)[0]\xrightarrow{\ \sim\ }
       R\Gamma_v(B_C\times\underline S,\Lambda_m).
\end{equation}
For the empty set both sides are zero.  Otherwise write $S=\varprojlim_i S_i$
over its finite quotients, including the one-point quotient.  The example
after \cite[Definition~7.8]{ScholzeDiamonds} constructs the affinoid
pro-\'etale limit $\operatorname{Spa}(C,\cO_C)\times S$ of the finite
unions indexed by $S_i$; the associated-diamond construction and
\cite[Example~11.12 and Lemma~15.2]{ScholzeDiamonds} identify its
diamond with $B_C\times\underline S$.
The point $\operatorname{Spa}(C,\cO_C)$ is strictly totally disconnected
by \cite[Definition~7.15 and Proposition~7.16]{ScholzeDiamonds}, so its
\'etale covers split and finite constant coefficients have no higher
cohomology there.  The preceding \cite[Lemma~15.6]{ScholzeDiamonds}
comparison transfers this assertion to the finite unions of their
diamonds.  The continuity statement
\cite[Proposition~14.9]{ScholzeDiamonds} computes the \'etale
cohomology on the limit as the filtered colimit of that on the finite
unions.  Its degree-zero group is
$\varinjlim_i C(S_i,\Lambda_m)=C_{\cts}(S,\Lambda_m)$ and its higher
groups vanish.  Proposition~14.10 of the same source, valid for the
coefficient ring $\Lambda_m$, transfers this calculation to the v-site.
Every comparison is induced by the constants map, and is natural in
$S$ and in automorphisms of $C$.

For compact $G$, apply \eqref{eq:geometric-point-constants} to
$S=G^a$ in every degree of \eqref{eq:finite-constant-cochains}.
It follows that $u_{B_C,G,m}$ is an equivalence.  The same argument
applies to the finite \'etale fiber of $(\Z/p^m)(t)$ over $C$,
where $p$ is invertible, retaining its arithmetic action.  The
compatible Tate line is finite free of rank one, so its tensor factor
commutes with the cochain terms and their derived coefficient limit.
Thus the integral comparison for
$B_CG$ is also compatible with Tate twists and arithmetic Galois
actions, without a choice of basis for the Tate fiber.  This
geometric-point calculation is used below for the compact groups
$P_n$; it is not asserted over $\Spd k$.

For descent with geometric coefficients it is useful to retain
profinite parameters.  Denote by $\mathscr C_{X,m}$ the condensed
derived complex whose derived value on a profinite set $S$ is
\[
 \mathscr C_{X,m}(S)=R\Gamma_v(X\times\underline S,\Lambda_m),
 \qquad \mathscr C_X=\Rlim_m\mathscr C_{X,m}.
\]
Descent for profinite covers gives these condensed complexes.  The value
of $\mathscr C_X$ at the point is $R\Gamma_{\mathrm{int}}(X)$.
A group acting on $X$ acts continuously
on this complex.  The coefficient complexes used below are
\emph{derived solid}: they lie in the derived category of solid
condensed modules.  Here is a concrete check of this condition.
On an affinoid perfectoid cover, the characteristic-$p$ structure
sheaf $\cO^\flat$ has no higher cohomology, and its values with
profinite parameters are continuous-function modules.  Their complete
linear topology expresses them as limits of discrete modules, so
they are solid.  Perfectoid hypercover descent computes the general
$\cO^\flat$ complex by limits of these objects; this is also the
construction in \cite[Theorem~3.2.1(iii) and Remark~3.2.2]{BMR2026}.
The Artin--Schreier
sequence gives the $\F_p$ complex, the coefficient exact sequences
give the $\Z/p^m$ complexes, and $\Rlim_m$ gives $\mathscr C_X$.
Derived solidity is preserved by these limits and extensions
\cite[Theorem~4.4(ii)]{Tang2026}.  This argument applies with the
group action and with all the profinite parameters retained.  It
also gives derived $p$-completeness of $\mathscr C_X$.

For compact $G$, the condensed standard resolution with its derived
parameter values retained gives
\begin{equation}\label{eq:geometric-group-descent}
 R\Gamma(G,\mathscr C_X)(*)\simeq
 \Rlim_m\operatorname{Tot}_{a}R\Gamma_v
        (X\times\underline{G^a},\Lambda_m)
 \simeq R\Gamma_{\mathrm{int}}([X/G]).
\end{equation}
Indeed all $G^a$ are profinite.  In the standard resolution, the
degree-$a$ derived Hom is precisely the derived parameter value on
$G^a$; this is the formal condensed construction in
\cite[\S2, pp.~5--6]{Anschutz2020}, before any acyclicity simplification.
The parameterized complex is bounded below, and the derived
totalization computes its point-valued invariants.  The second
comparison is quotient \v{C}ech descent at each finite level; the
derived limit and totalization commute because both are limits.
This is the geometric descent comparison used here.  On the same bar
terms, using \cite[Lemma~2.2]{Anschutz2020} for the source parameter
value, the finite unit $\underline{\Lambda_m}\to\mathscr C_{X,m}$
is exactly the map in \eqref{eq:finite-constant-cochains}.  The complexes
$\mathscr C_{X,m}$ and $\mathscr C_X$ are concentrated in cohomological
degrees at least zero, so their units factor through their degree-zero
cohomology by the truncation triangle.  Thus the derived-invariants map
agrees with $u_{X,G,\mathrm{int}}$ after the coefficient limit, and its
factorization through degree zero gives the constants bottom-row edge map.

The next lemma makes these descent groups computable from a finite
complex. Its final assertion will let us use arithmetic weights known
on ordinary cohomology groups while retaining the solid coefficients.

\begin{lemma}[Finite resolutions and scalar actions]
\label{lem:solid-resolutions}
Let $G$ be a compact $p$-adic analytic group without elements of order
$p$, and let $\Lambda_G=\Z_p[[G]]$.  There is a finite resolution
$P_\bullet\to\Z_p$ by finitely generated projective profinite
$\Lambda_G$-modules.  Its length can be taken to be $\dim G$.
For any derived-solid $\Z_p$-complex $K$ with a continuous $G$-action,
the point-valued derived invariants are computed by
\begin{equation}
\label{eq:finite-resolution}
 R\Gamma(G,K)(*)\simeq
 \Hom_{\Lambda_G}(P_\bullet,K).
\end{equation}
Here the right side means the finite total complex of morphisms in
solid modules.  In particular, invariants of a module in the heart
have cohomology only in degrees $0,\ldots,\dim G$.

Suppose $K$ is bounded below and an operator $\gamma$ commutes with
$G$.  If $\gamma$ acts by $\lambda_j\in\Q_p$ on $H^j(K)(*)[1/p]$,
then it acts by the same scalar on every term in the $j$th row of
the rationalized descent spectral sequence
\[
 H^s\bigl(G,H^j(K)\bigr)(*)[1/p]
       \ \Longrightarrow\ H^{s+j}\bigl(R\Gamma(G,K)(*)\bigr)[1/p].
\]
For $K=\mathscr C_X$, the abutment is $H^{s+j}_{\mathrm b}([X/G])$.
\end{lemma}

\begin{proof}
First construct the resolution in profinite modules. A compact $p$-adic analytic group has a noetherian completed group
ring; if it has no $p$-torsion, then
$\operatorname{pd}_{\Lambda_G}\Z_p=\operatorname{cd}_p(G)=\dim G$.
Here projective dimension is in profinite modules.
These results for compact groups, including the agreement with
abstract finitely generated modules, are recorded in
\cite[\S\S1.1--1.2, pp.~275--276]{Venjakob2002}.
The noetherianity input for analytic pro-$p$ groups is
\cite[V.2.2.4]{Lazard1965}; the compact-group statement uses the
finite-index passage recorded by Venjakob.
Choose successive finite free presentations.  Noetherianity keeps
their kernels finitely generated, and the projective-dimension
bound makes the last syzygy projective.  Their natural compact
topologies give the asserted profinite resolution.

To apply the resolution to the geometric coefficient complexes, we
pass from continuous actions to solid modules over the completed group
ring. Write
$A=\underline{\Z_p}[\underline G]$ for the free condensed group
ring.  A continuous action on a complex means an object of
$D(A\text{-}\mathrm{Cond})$ with derived-solid underlying complex.
By \cite[Lemma~1.3]{Anschutz2020}, the group-ring analytic
structure has underived solidification; its normalized ring is
\[
 A^\square\simeq\underline{\Z_p[[G]]}
   =\underline{\varprojlim_{m,V}(\Z/p^m)[G/V]},
\]
where $V$ runs through open normal subgroups.  The ring
identification, including multiplication, is
\cite[\S3.6, preceding Proposition~3.21]{Tang2026}.
Analyticity gives fully faithful embeddings of the \emph{unbounded}
derived category
of solid $G$-modules into both $D(A\text{-}\mathrm{Cond})$ and
$D(A^\square\text{-}\mathrm{Cond})$, as stated in
\cite[\S1, pp.~3--4]{Anschutz2020}.
The solid heart is characterized by solidity of the underlying
condensed module. The essential-image criterion of Clausen--Scholze's
analytic theory extends this characterization to unbounded complexes:
membership is detected on cohomology objects
\cite[Proposition~12.4 and Remark~12.5]{ScholzeAnalyticGeometry}.
Thus its objects are precisely the complexes with derived-solid
underlying coefficients. Consequently the external derived Hom
complexes satisfy
\[
 R\Gamma(G,K)(*)
 \simeq\RHom_{A\text{-}\mathrm{Cond}}(\underline{\Z_p},K)
 \simeq\RHom_{G\text{-}\mathrm{Solid}}(\underline{\Z_p},K).
\]
The first identification is the external derived Hom for the condensed
classifying topos $BG_{\mathrm{pro\acute et}}$
\cite[\S2, pp.~5 and~8]{Anschutz2020}; the last category is that of
solid $\Z_p[[G]]$-modules by
\cite[Proposition~3.21]{Tang2026}.
They apply to the given complex $K$ itself, including when its
coefficients are neither profinite nor derived $p$-complete.

The profinite resolution stays exact and projective in solid
modules \cite[Theorems~3.2 and~3.14(i)]{Tang2026}.
Exactness can be seen by testing on extremally
disconnected profinite sets: maps from such a set lift through a
surjection of profinite spaces.  Each $P_i$ is a retract of
$\Lambda_G^{r_i}$, and
\[
 \Hom_{\Lambda_G}(\Lambda_G,M)=M(*).
\]
Evaluation at the point is exact, since the point is extremally
disconnected.  This proves projectivity against an arbitrary solid
coefficient module $M$, even when $M$ is not profinite.  A bounded
projective resolution computes derived Hom against any complex,
which proves the displayed calculation and the amplitude assertion.

Finally, $\Hom_{\Lambda_G}(P_i,H^j(K))$ is a retract of a finite
sum of $H^j(K)(*)$.  The retract commutes with $\gamma$, because
$\gamma$ commutes with the $G$-action.  After inverting $p$, the
entire finite complex computing the $j$th row therefore has scalar
action $\lambda_j$.  Its cohomology has the same action.
This uses a scalar identity on point values only; it does not
deduce an identity of condensed sheaves from their point values.
Boundedness of the resolution gives a finite filtration in each
total degree, so exact rationalization preserves the spectral
sequence and its abutment.
\end{proof}

\subsection{Compact groups and the affine building}

Put
\[
 J_n=\{g\in\GL_n(\Q_p):|\det g|_p=1\}\qquad(n=2,3).
\]
Let $D_n$ denote the division algebra of degree $n$ whose maximal
order has unit group $P_n$.  We need only the following low-degree
part of the rational Lie algebra calculation.

\begin{proposition}[The group classes]
\label{prop:group-classes}
For $G=P_n$ or $J_n$, where $n=2,3$,
\[
 H^1_{\mathrm b}(G)=\Q_p,
 \qquad H^2_{\mathrm b}(G)=0,
 \qquad H^3_{\mathrm b}(G)=\Q_p.
\]
For $n=2,3$ and every compact open subgroup $K\subset J_n$,
restriction gives an isomorphism in every degree, with the compact
comparison identifying its target:
\[
 H^*_{\mathrm b}(J_n)\xrightarrow{\mathrm{res}}H^*_{\mathrm b}(K)
 \cong H^*(\mathfrak{gl}_n(\Q_p),\Q_p).
\]
In degree three the block inclusion $J_2\hookrightarrow J_3$
induces a nonzero map.  Fix arbitrary nonzero classes
\[
 e_{P_n}\in H^3_{\mathrm b}(P_n),\qquad
 e_{J_n}\in H^3_{\mathrm b}(J_n).
\]
\end{proposition}

\begin{proof}
For an inclusion \(K_1\subset K_2\) of compact open subgroups of
\(\GL_n(\Q_p)\), choose a \(K_2\)-stable lattice \(\Lambda\) and then
\(r\) large enough that
\[
 U=1+p^r\End_{\Z_p}(\Lambda)\subset K_1.
\]
This is a uniform subgroup normal in \(K_2\). In the division-algebra
case one uses a sufficiently deep congruence subgroup of
\(\cO_{D_n}^{\times}\), which is normal in the relevant compact group.
For these small groups, Lazard's comparison is
\cite[Theorem~3.3.3 and Remark~3.3.4]{HKN2011} with trivial
\(\Z_p\)-coefficients, followed by inversion of \(p\).
Its coefficient naturality, rational agreement, and group naturality are
\cite[Theorem~3.1.1(1)--(3)]{HKN2011}, with numbering in the arXiv
revision specified in the bibliography. Equip the uniform groups with
their lower \(p\)-series. At these odd primes the proof of
Theorem~3.3.3 gives \(e=1\), and the associated graded target is generated
in degree one, as required by the group-naturality clause. Inclusions and
the block homomorphisms preserve these filtrations. The common subgroup
\(U\) therefore identifies restriction with the identity on the same
ambient Lie algebra cohomology.

Here is why passing back to a compact group is valid even for an index
divisible by \(p\). For \(U\triangleleft K\), the augmentation and norm
of the finite permutation module \(\Z[K/U]\) induce, by continuous
Shapiro, the identities
\[
 \operatorname{cor}\operatorname{res}=[K:U],\qquad
 \operatorname{res}\operatorname{cor}
      =\sum_{q\in K/U}q^*
\]
at each finite coefficient level, compatibly with reduction in \(m\).
Shapiro here is the derived induction--restriction adjunction: induction
for an open subgroup is a direct sum on underlying condensed modules and
is exact by (AB4) \cite[Theorems~2.2 and~2.4]{Tang2026}; restriction is
also exact. Their adjunction therefore descends directly to derived
categories. Its naturality identifies the
permutation maps with the actual restrictions and conjugations in the
finite continuous bar complexes. Apply \(\Rlim_m\) to these identities
before taking cohomology and inverting \(p\). Only then divide by
\([K:U]\). This gives
\(H^*_{\mathrm b}(K)\xrightarrow{\sim}H^*_{\mathrm b}(U)^{K/U}\)
by restriction, with the compact
continuous-cochain interpretation of \eqref{eq:group-cochain-bridge}.

The finite quotient acts trivially in every Lie degree. The
Chevalley--Eilenberg complex is a finite-dimensional algebraic
\(\GL_n\)-representation after a splitting-field extension. Cartan's
formula \(\mathcal L_X=d\iota_X+\iota_Xd\) makes the derivative of its
action on cohomology zero. In characteristic zero, an algebraic action
of the connected group \(\GL_n\) with zero derivative is trivial.
Thus all inner automorphisms act trivially, in every degree; the same
conclusion descends for the division algebra. This proves the compact
restriction identification without an all-degree numerical calculation.

After a finite splitting-field extension, either Lie algebra is
$\mathfrak{gl}_n$.  The finite Chevalley--Eilenberg calculation in
\Cref{app:background} gives its asserted groups in degrees one through
three, with degree-three representative
\[
 (X,Y,Z)\longmapsto\tr\bigl(X[Y,Z]\bigr),
\]
or its reduced-trace version for $D_n$.
This is the invariant-form cocycle of
\cite[\S21]{ChevalleyEilenberg1948} and \cite[\S11]{Koszul1950};
the appendix records the normalization and restriction used here.
These representatives are defined over $\Q_p$ and invariant under
inner automorphisms, so the assertions descend from the splitting
field.  Restriction of the displayed form to the upper-left or
lower-right $2\times2$ block is the corresponding form for
$\mathfrak{gl}_2$, whose class is nonzero.  This proves the compact
calculations and their compatibility.

To pass to $J_n$, let $\mathcal B_n$ be the reduced affine building
of $\GL_n(\Q_p)$.  It is contractible of dimension $n-1$.
The group $J_n$ preserves vertex types, its face stabilizers
$K_\sigma$ are compact open, and a chamber $\Delta$ is a
fundamental domain, including its faces.  A proof of these building
properties by lattice norms is included in \Cref{app:background}.
Write its bounded augmented cellular resolution as
\[
 P_r=\bigoplus_{\substack{\sigma\subset\Delta\\\dim\sigma=r}}
       \Z[\underline{J_n/K_\sigma}].
\]
It is exact also as a complex of condensed modules. Indeed, on a
profinite parameter set a continuous section of a discrete cellular
chain group has finite image. Contractibility lifts each of the finitely
many cycles in that image; choosing those lifts on its clopen fibers
gives a continuous lift. The same exact derived induction--restriction
adjunction therefore supplies continuous Shapiro for each summand.
At coefficient level \(\Lambda_m\), applying derived Hom gives a finite
filtration with graded terms
\(\bigoplus_{\dim\sigma=r}C_{\cts}^{\bullet}(K_\sigma,\Lambda_m)[-r]\)
in \(D(\mathrm{Ab})\).
A cellular boundary is the literal signed projection of the corresponding
coset modules, so naturality of that adjunction makes each face map the
signed restriction between stabilizers, with no index multiplier.
All of these statements are compatible with coefficient reduction.

Take \(\Rlim_m\) of this finite filtration. There are finitely many
orbit terms in each of finitely many cellular degrees, so this step uses
only finite sums and cones in the cellular direction. Taking cohomology of
the resulting finite filtration gives the integral spectral sequence
\[
 E_1^{r,s}=
 \bigoplus_{\substack{\sigma\subset\Delta\\\dim\sigma=r}}
 H^s_{\mathrm{int}}(K_\sigma)
 \ \Longrightarrow\ H^{r+s}_{\mathrm{int}}(J_n),
\]
with precisely those signed restriction faces. It does not interchange
an infinite group-cochain totalization with the derived inverse limit.
Only after this construction do we invert \(p\); exact rationalization
preserves the finite filtration and its abutment.

Under the compact comparison, every restriction between face
stabilizers is the identity on the same Lie algebra cohomology;
any change of representatives contributes an inner conjugation,
which also acts trivially.  Thus each rationalized row is the
cellular cochain complex of the simplex $\Delta$ with constant
coefficients $H^s(\mathfrak{gl}_n,\Q_p)$.
It has cohomology only in cellular degree zero.
The resulting edge map is restriction to a vertex stabilizer.
Restriction further to compact open subgroups, using intersections
and the compact comparison, proves the asserted canonical
identifications.  The block restriction can now be tested on
compact open subgroups, where it is the trace calculation above.
\end{proof}

\subsection{The parabolic restriction}

The relevant subgroup remembers a lattice only in each determinant
line.  In a basis adapted to a line it is
\[
 D=\left\{
 \begin{pmatrix}a&v\\0&A\end{pmatrix}:
 a\in\Z_p^\times,\quad A\in J_2,\quad
 v\in\Q_p^{1\times2}\right\}\subset J_3.
\]
Write $\rho_3:D\hookrightarrow J_3$ for the inclusion and
$\rho_2:D\to J_2$ for the quotient action.

\begin{lemma}[Unipotent acyclicity and restriction]
\label{lem:parabolic}
Inflation induces an isomorphism
\[
 H^*_{\mathrm b}(\Z_p^\times\times J_2)
       \xrightarrow{\ \sim\ }H^*_{\mathrm b}(D).
\]
There is a scalar $b\in\Q_p^\times$ such that
\[
 \rho_3^*e_{J_3}=b\,\rho_2^*e_{J_2}.
\]
\end{lemma}

\begin{proof}
The kernel of $D\to\Z_p^\times\times J_2$ is
$N=\Q_p^2$, with its additive topology.
Fix $m$ and exhaust it by the compact open subgroups
$N_r=p^{-r}\Z_p^2$ for $r\geq0$.  Their finite-coefficient
cohomology, computed by the two-variable Koszul resolution, is
\[
 H^q(N_r,\Z/p^m)=\bigwedge^q(\Z/p^m)^2.
\]
The standard bases identify restriction from $N_{r+1}$ to $N_r$
with multiplication by $p^q$ in degree $q$: restriction multiplies
each degree-one homomorphism by $p$, and is compatible with cup
products.  For $q>0$ this inverse system is pro-zero; for $q=0$
it is constant with identity transition maps.

This calculation computes the cohomology of $N$ without a missing
derived-limit term.  Indeed a continuous cochain on $N$ is exactly
a compatible family of cochains on the $N_r$.
Restriction of cochains is surjective, since $N_r^a$ is clopen in
$N_{r+1}^a$ and a function extends by zero on its complement.
Thus the ordinary inverse limit of the cochain complexes realizes
their derived inverse limit.  The Milnor exact sequence has
$\varprojlim^1_r=0$ in every degree: this holds both for a constant
tower with identity maps and for a pro-zero tower.  It follows that
the constant map is an equivalence
\[
 \Lambda_m\xrightarrow{\ \sim\ }C_{\cts}^{\bullet}(N,\Lambda_m).
\]
The same proof works with coefficients
$C_{\cts}(S,\Z/p^m)$ for every profinite parameter space $S$.
The finite Koszul complex gives finite sums of this module, the
positive-degree restrictions are again $p^q$, and extension by
zero is continuous on $N_r^a\times S$.
The condensed standard resolution evaluated on each such parameter
identifies this calculation with the internal equivalence
\[
 \underline{\Lambda_m}\xrightarrow{\ \sim\ }
       R\Gamma(N,\underline{\Lambda_m}).
\]
Thus it proves an equivalence of coefficient objects for extension
descent.  The equivalence is the map of constants, so it is equivariant
under the Levi action.

Now take $\Rlim_m$.  In degree zero the parameter modules have
surjective transition maps, by lifting locally constant values,
and their limit is $C_{\cts}(S,\Z_p)$.
Hence the internal coefficient limit satisfies
\[
 \underline{\Z_p}\xrightarrow{\ \sim\ }
       \Rlim_m R\Gamma(N,\underline{\Lambda_m}).
\]
At each finite level, internal descent for the split extension identifies
$R\Gamma(D,\underline{\Lambda_m})$ with
$R\Gamma(\Z_p^\times\times J_2,
R\Gamma(N,\underline{\Lambda_m}))$.
The Levi-equivariant constant equivalence therefore induces inflation.
After point evaluation, the standard-resolution comparison identifies
this map with finite continuous-cochain inflation.  Taking $\Rlim_m$
and using \eqref{eq:integral-group-cochains} gives the claimed
isomorphism integrally before rationalization.

Finally, the compact resolution for $\Z_p^\times$ and the finite
building resolution for $J_2$ give the usual rational product
calculation.  Since $H^*_{\mathrm b}(\Z_p^\times)$ is exterior
on one degree-one generator and $H^2_{\mathrm b}(J_2)=0$,
degree three of the Levi is exactly $H^3_{\mathrm b}(J_2)$.
Thus $\rho_3^*e_{J_3}$ is a scalar multiple of
$\rho_2^*e_{J_2}$.  Restricting to the section
$A\mapsto\operatorname{diag}(1,A)$ shows that this scalar is
nonzero, by \cref{prop:group-classes}.
\end{proof}

The vanishing just proved is specific to our order of operations.
For example, continuous homomorphisms $\Q_p\to\Q_p$ give a
nonzero $H^1_{\cts}(\Q_p,\Q_p)$, whereas the preceding integral
calculation gives $H^1_{\mathrm b}(\Q_p)=0$.

\subsection{The integral Drinfeld theorem}

Let $\Omega_C^d$ be the base change to $C$ of Drinfeld's symmetric
space over $\Q_p$, the complement of the $\Q_p$-rational
hyperplanes in $\mathbf P^d$.  The following precise input supplies
the weights used in \cref{prop:primitive-comparison}.

\begin{theorem}[Colmez--Dospinescu--Nizio\l]
\label{thm:drinfeld}
For $0\leq j\leq d$ there are compatible
$\GL_{d+1}(\Q_p)\times\Gal(\overline\Q_p/\Q_p)$-equivariant
isomorphisms of groups
\[
 H^j_{\mathrm{int}}(\Omega_C^d)(j)
       \simeq\operatorname{Sp}_j(\Z_p)^*,
 \qquad
 H^j_{\mathrm{\acute et}}(\Omega_C^d,\F_p(j))
       \simeq\operatorname{Sp}_j(\F_p)^*.
\]
Here $\operatorname{Sp}_j$ is the generalized Steinberg
representation of \cite[\S4.1]{CDN2021}, the star denotes its
continuous coefficient-linear dual with weak topology, and
the right sides have trivial Galois action.  These groups vanish
for $j>d$, and $H^0_{\mathrm{int}}(\Omega_C^d)=\Z_p$ by constants.
In particular, on the untwisted integral group in degree $j$,
an arithmetic Galois element $\gamma$ acts by
$\chi_{\mathrm{cyc}}(\gamma)^{-j}$.
\end{theorem}

The group calculation is \cite[Theorem~1.1]{CDN2021}, with its integral
cohomology convention.  The comparison with
$H^j_{\mathrm{pro\acute et}}(\Omega_C^d,\widehat\Z_p(j))$
is made explicitly in the proof in \cite[\S5]{CDN2021}; finite
coefficient site comparison and $\Rlim_m$ identify it with the
convention above.  Thus the scalar assertion holds on the
integral point groups themselves.  To identify the stated degree-zero
map, choose a $C$-point of $\Omega_C^d$ by taking $d+1$ coordinates
linearly independent over $\Q_p$.  By
\eqref{eq:geometric-point-constants}, evaluation there splits the
constants unit at each finite level and after $\Rlim_m$.  The cited
calculation makes $H^0_{\mathrm{int}}(\Omega_C^d)$ a free rank-one
$\Z_p$-module, so this split inclusion from $\Z_p$ is an isomorphism.
This identifies the actual constants map locally from the degree-zero
calculation.  Applying
\cref{lem:solid-resolutions} to a compact open uniform subgroup
of $J_{d+1}$ propagates it to the group-cohomology rows, without
any assumption about cohomology of a rationalized sheaf.
